\chapter{进程}

\section{进程引入}
\concept{前趋图}
\begin{points}
  \item 前趋图是有向无循环图，用于描述程序段、语句或进程之间的先后执行关系。
  \item 若 $P_i \rightarrow P_j$，表示 $P_i$ 必须在 $P_j$ 开始之前完成。
  \item 没有前趋的结点称为初始结点，没有后继的结点称为终止结点。
  \item 做题时根据箭头判断哪些操作必须先完成，哪些操作可以并发。
\end{points}
\plain{用有向无环图画出“谁必须在谁之前完成”，目的是识别哪些操作能并发。}
\exam{常考给前趋关系画图、数可并发的语句、写同步信号量约束。}


\concept{程序的顺序执行}
\begin{points}
  \item 顺序执行指一个程序中的多个操作严格按照规定顺序执行，例如输入 $I \rightarrow$ 计算 $C \rightarrow$ 输出 $P$。
  \item 特征一：顺序性，每一步必须在下一步开始前完成。
  \item 特征二：封闭性，程序运行结果不受外界因素影响。
  \item 特征三：可再现性，只要初始条件和执行环境相同，每次运行结果相同。
\end{points}
\plain{一个程序自己按固定步骤走，外界不插手，所以结果可重复。}
\exam{常考顺序性、封闭性、可再现性三个特征。}


\concept{程序的并发执行}
\begin{points}
  \item 并发执行指多个程序或程序段在时间上重叠推进，一个程序尚未结束，另一个程序已经开始。
  \item 特征一：间断性，程序表现为“执行 - 暂停 - 执行”。
  \item 特征二：失去封闭性，多个程序共享资源，资源状态可能被他人改变。
  \item 特征三：不可再现性，在初始条件相同的情况下，执行顺序不同也可能导致结果不同。
\end{points}
\plain{核心是“看起来同时推进”。单 CPU 上靠时间片交替，多 CPU 上才可能真的同一时刻运行。}
\exam{常考并发与并行区别、单处理机下最多几个进程运行、并发导致为什么需要同步互斥。}

\section{进程定义与描述}
\concept{进程的定义}
\begin{points}
  \item 进程是程序的一次执行过程，是系统进行资源分配和调度的基本单位之一。
  \item 进程不仅包含程序代码，还包含数据、堆栈、打开文件、寄存器现场、程序计数器和 PCB 等运行信息。
  \item 进程具有动态性、并发性、独立性、异步性和结构性。
\end{points}
\plain{程序是静态代码，进程是这段代码跑起来后的动态活动，还带资源和现场。}
\exam{常考程序与进程区别、进程实体由程序段/数据段/PCB 构成、进程是资源分配基本单位。}


\concept{程序与进程的区别}
\begin{points}
  \item 程序是静态的代码文件，进程是动态的执行活动。
  \item 一个程序可以对应多个进程，例如多个用户同时打开同一个编辑器。
  \item 一个进程在运行中可能执行多个程序片段，例如通过 \code{exec} 装入新程序。
  \item 程序没有生命周期状态，进程有创建、就绪、运行、阻塞、终止等状态。
\end{points}
\plain{程序是静态代码，进程是这段代码跑起来后的动态活动，还带资源和现场。}
\exam{常考程序与进程区别、进程实体由程序段/数据段/PCB 构成、进程是资源分配基本单位。}


\concept{PCB 进程控制块}
\begin{points}
  \item PCB 是操作系统感知和管理进程的唯一标志。
  \item PCB 保存进程标识信息、处理机状态、调度信息、进程控制信息、内存管理信息、I/O 状态信息等。
  \item 上下文切换本质上是保存当前进程 PCB 中的现场，并从下一个进程 PCB 中恢复现场。
  \item 选择题直觉：问“进程存在的唯一标志”通常答 PCB，不是程序、PID 或页表。
\end{points}
\plain{PCB 就是 OS 给每个进程建的档案，没有 PCB，OS 就没法调度和恢复这个进程。}
\exam{常考 PCB 内容、PCB 的作用、进程存在的唯一标志。}


\section{进程状态与转换}
\concept{三状态模型}
\begin{points}
  \item \term{就绪态}：进程已具备运行条件，只差 CPU。
  \item \term{运行态}：进程正在 CPU 上执行。
  \item \term{阻塞态/等待态}：进程因等待某事件不能运行，即使 CPU 空闲也不能执行。
  \item 状态转换：就绪 $\rightarrow$ 运行由调度引起；运行 $\rightarrow$ 就绪由时间片到或被抢占引起；运行 $\rightarrow$ 阻塞由等待 I/O/资源引起；阻塞 $\rightarrow$ 就绪由等待事件完成引起。
\end{points}
\plain{进程不是一直运行，它会在就绪、运行、阻塞等状态之间因调度或等待事件而移动。}
\exam{常考三/五/七状态转换图，尤其运行到阻塞、阻塞到就绪、就绪到运行的触发条件。}


\concept{五状态模型}
\begin{points}
  \item 在三状态基础上加入新建态和终止态。
  \item 新建态表示进程正在创建，尚未进入就绪队列。
  \item 终止态表示进程执行结束或被撤销，等待系统回收资源。
  \item 创建进程通常涉及分配 PCB、建立地址空间、装入程序、初始化资源、进入就绪队列。
\end{points}
\plain{进程不是一直运行，它会在就绪、运行、阻塞等状态之间因调度或等待事件而移动。}
\exam{常考三/五/七状态转换图，尤其运行到阻塞、阻塞到就绪、就绪到运行的触发条件。}


\concept{挂起状态}
\begin{points}
  \item 当内存紧张或系统需要调节多道程序度时，可把进程换出到外存，形成挂起状态。
  \item 就绪挂起：进程不在内存，但一旦调入内存即可运行。
  \item 阻塞挂起：进程不在内存，且还在等待某事件。
  \item 中级调度/交换调度常与挂起状态相关。
\end{points}
\plain{进程不是一直运行，它会在就绪、运行、阻塞等状态之间因调度或等待事件而移动。}
\exam{常考三/五/七状态转换图，尤其运行到阻塞、阻塞到就绪、就绪到运行的触发条件。}


\section{进程组织、控制与管理}
\concept{进程队列}
\begin{points}
  \item 就绪队列：所有驻留内存且处于就绪状态、等待 CPU 的进程集合。
  \item 设备队列：等待某个 I/O 设备的进程集合。
  \item 作业队列：外存上等待进入内存的作业集合。
  \item 进程在生命周期中会在多个队列之间迁移，调度程序就是从这些队列中选择合适对象。
\end{points}
\plain{队列是 OS 组织进程的方式：就绪的排就绪队列，等设备的排设备队列。}
\exam{常考进程在生命周期中如何在不同队列迁移。}


\concept{进程创建与终止}
\begin{points}
  \item 创建进程需要申请唯一标识、分配 PCB、分配必要资源、初始化上下文、插入就绪队列。
  \item 进程终止后，系统要回收内存、打开文件、I/O 资源、PCB 等。
  \item Unix/Linux 中常见系统调用：\code{fork()} 创建子进程，\code{exec()} 装入新程序，\code{wait()} 等待子进程，\code{exit()} 退出进程。
  \item \code{fork()} 返回两次：父进程中返回子进程 PID，子进程中返回 0；失败返回负值。
\end{points}
\plain{进程控制就是 OS 创建/撤销/阻塞/唤醒/切换进程，都会改 PCB 和队列。}
\exam{常考 fork/exec/wait/exit 以及父子进程数量。}


\concept{进程切换与模式切换}
\begin{points}
  \item 模式切换是用户态和内核态之间的切换，不一定更换正在运行的进程。
  \item 进程切换是 CPU 从一个进程切换到另一个进程，一定需要内核介入，通常伴随上下文保存和恢复。
  \item 系统调用会导致模式切换；是否导致进程切换，要看是否阻塞或调度。
  \item I/O 请求通常先从用户态陷入内核态，若需要等待 I/O 完成，则当前进程阻塞并发生进程切换。
\end{points}
\plain{模式切换只是换权限；进程切换还要保存/恢复另一个进程的现场，开销更大。}
\exam{常考系统调用是否必然进程切换、进程切换需要保存哪些上下文。}


\section{进程通信}
\concept{低级通信与高级通信}
\begin{points}
  \item 低级通信主要通过信号量等同步工具传递少量控制信息。
  \item 高级通信用于传递大量数据或结构化消息，常见方式包括共享存储、消息传递、管道。
\end{points}
\plain{进程地址空间相互隔离，所以要么共享一块区域，要么通过内核传消息。}
\exam{常考共享存储、消息传递、管道的特点和适用场景。}


\concept{共享存储}
\begin{points}
  \item 多个进程共享同一块内存区域，通过读写共享区交换数据。
  \item 优点：速度快，适合大量数据交换。
  \item 缺点：必须另行解决同步互斥问题，防止同时读写导致不一致。
\end{points}
\plain{多个进程都想用同一批资源，所以必须规定能不能一起用、谁先用、用完怎么释放。}
\exam{常考互斥共享/同时共享举例，或者问并发和共享为什么是 OS 最基本特征。}


\concept{消息传递}
\begin{points}
  \item 进程通过发送和接收消息通信，可由操作系统提供消息队列、邮箱等机制。
  \item 直接通信：发送者明确指定接收者。
  \item 间接通信：通过邮箱或消息队列进行。
  \item 优点：同步控制相对清晰，适合分布式系统；缺点是消息复制和内核介入可能带来开销。
\end{points}
\plain{进程地址空间相互隔离，所以要么共享一块区域，要么通过内核传消息。}
\exam{常考共享存储、消息传递、管道的特点和适用场景。}


\concept{管道}
\begin{points}
  \item 管道是一种半双工通信机制，常用于具有亲缘关系的进程之间。
  \item 匿名管道常用于父子进程；命名管道可用于无亲缘关系的进程。
  \item Shell 中的 \code{|} 就是把前一个命令的输出作为后一个命令的输入。
\end{points}
\plain{进程地址空间相互隔离，所以要么共享一块区域，要么通过内核传消息。}
\exam{常考共享存储、消息传递、管道的特点和适用场景。}


